\section{Replacing a module}\label{section: replacement}

\begin{lemma}[Seams]\label{lemma: seam}
  For every star-free term $G$ there is a finite set $\mathcal{R}(G)$ of pairs $(A, C)$ of star-free terms with the following properties.
  We have $\cd(A), \cd(C) \le \cd(G)$ for all pairs.
  Moreover, let $\A$ be a valuation and $x \le c \le y$ such that no input pair $(a, a')$ satisfies $x \le a < c < a' \le y$.
  Then
  \[
    G(x, y) \iff \bigvee_{(A, C) \in \mathcal{R}(G)} A(x, c) \wedge C(c, y).
  \]
\end{lemma}
\begin{proof}
  By induction on $G$.
  For a variable $o \in \set{\vr, \vc, \vb}$ we take $\set{(\id, o), (o, \id)}$.
  If $x < c < y$, then $o(x, y)$ fails by assumption, and if $c \in \set{x, y}$, one of the two pairs applies.
  For $\id$ we take $\set{(\id, \id)}$, and for $\emp$ the empty set.
  For Boolean operations, let $\mathcal{A}$ and $\mathcal{C}$ be the sets of first and second components occurring in the pairs of the subterms.
  The atoms of the finite Boolean algebras generated by $\mathcal{A}$ and by $\mathcal{C}$ are Boolean combinations of terms of depth at most $\cd(G)$.
  Every Boolean combination of the subterms is a union of products of such atoms.
  For complement this uses that $\topr(x, y)$ holds for $x \le y$.
  For $G = G_1 \compo G_2$, split on whether the intermediate point $z$ satisfies $z \le c$ or $z \ge c$.
  If $z \le c$, the interval $\range{z, y}$ again satisfies the hypothesis, so $G_2(z, y)$ decomposes at $c$.
  This gives the pairs $(G_1 \compo A_2, C_2)$ for $(A_2, C_2) \in \mathcal{R}(G_2)$.
  Symmetrically, $z \ge c$ gives $(A_1, C_1 \compo G_2)$ for $(A_1, C_1) \in \mathcal{R}(G_1)$.
  In both cases the depth is at most $\max(\cd(G_1), \cd(G_2)) + 1$.
\end{proof}

\begin{lemma}[Replacement]\label{lemma: replacement}
  Let $\A, \A'$ be valuations on the same linear order $V$, and let $u < v$ in $V$.
  Let $O \defeq V \setminus \set{z \mid u < z < v}$ be the set of \emph{outer} vertices.
  Assume the following.
  \begin{enumerate}
    \item\label{item: same outside} $\A$ and $\A'$ have the same input pairs inside $O \times O$.
    \item\label{item: isolated} In neither valuation does an input pair have one component in $\set{z \mid u < z < v}$ and the other in $O \setminus \set{u, v}$.
    \item\label{item: local equivalence} $\A\range{u, v} \eqq{q} \A'\range{u, v}$.
  \end{enumerate}
  Then for all $F \in \SF{q}$ and $x, y \in O$: $F(x, y)$ holds in $\A$ iff it holds in $\A'$, and the same holds for $F^{*}(x, y)$.
\end{lemma}
\begin{proof}
  \emph{Star-free terms.}
  We use induction on $F \in \SF{q}$.
  Atoms are handled by (\ref{item: same outside}), and Boolean operations are immediate.
  Let $F = G_1 \compo G_2$ with $\cd(G_1), \cd(G_2) \le q - 1$, and let $x \le z \le y$.
  If $z \in O$, we use the induction hypothesis for $G_1(x, z)$ and $G_2(z, y)$.
  Otherwise $x \le u < z < v \le y$.
  By (\ref{item: isolated}), \Cref{lemma: seam} applies to $G_1(x, z)$ at the seam $u$ and to $G_2(z, y)$ at the seam $v$.
  So the existence of such a $z$ is equivalent to
  \[
    \bigvee_{(A_1, C_1) \in \mathcal{R}(G_1)}\ \bigvee_{(A_2, C_2) \in \mathcal{R}(G_2)}
      A_1(x, u) \wedge \bigl((C_1 \inter \Dt) \compo (A_2 \inter \Dt)\bigr)(u, v) \wedge C_2(v, y).
  \]
  The intervals $\range{x, u}$ and $\range{v, y}$ consist of outer vertices.
  By (\ref{item: same outside}) and \Cref{lemma: locality}, the values of $A_1(x, u)$ and $C_2(v, y)$ agree in $\A$ and $\A'$.
  The middle conjunct is evaluated in $\range{u, v}$ by a term of depth at most $q$.
  So it agrees by \Cref{lemma: spectrum}.

  \emph{Stars.}
  Let $F^{*}(x, y)$ hold in $\A$, witnessed by $x = x_0 < x_1 < \dots < x_k = y$ with $F(x_{i-1}, x_i)$.
  We may drop repeated points, since all values are forward.
  If no $x_i$ is in the interior, the path transfers to $\A'$ by the first part.
  Otherwise the interior points form a block $x_s, \dots, x_t$ with $1 \le s \le t \le k - 1$.
  Let $a \defeq x_{s-1} \le u$ and $b \defeq x_{t+1} \ge v$.
  By \Cref{lemma: seam} at $u$ and $v$, there are $(A_1, C_1), (A_2, C_2) \in \mathcal{R}(F)$ with the following properties.
  $A_1(a, u)$ and $C_2(v, b)$ hold, and so does the observation
  \[
    \bigl((C_1 \inter \Dt) \compo F^{*} \compo (A_2 \inter \Dt)\bigr)(u, v).
  \]
  These terms lie in $\SF{q}$.
  By \Cref{lemma: spectrum} the observation also holds in $\A'\range{u, v}$, which gives interior $z \le z'$ there.
  By \Cref{lemma: seam} in $\A'$, we get $F(a, z)$ and $F(z', b)$.
  Here $A_1(a, u)$ and $C_2(v, b)$ are unchanged by (\ref{item: same outside}).
  The parts $x_0, \dots, x_{s-1}$ and $x_{t+1}, \dots, x_k$ consist of outer vertices and transfer by the first part.
  So $F^{*}(x, y)$ holds in $\A'$.
  The converse direction is symmetric.
\end{proof}
